\subsection*{Motivation}

Chapters 29 and 30 established the MDP machinery and the maximum-entropy framework; Chapter 28 derived the closed-form DPO solution to the KL-regularized RLHF objective. What remains is the \emph{online} alignment loop: an algorithm that takes a pre-trained language model (Chapter 27), a reward signal, and produces a better policy by gradient ascent on $\mathbb{E}_{y\sim\pi_\theta}[r(x,y)]$. We derive that loop from first principles, starting with the policy gradient theorem (Sutton, 1999), reducing variance with baselines and GAE, enforcing a trust region with PPO, dropping the value function with GRPO (DeepSeek, 2024), and ending by running GRPO on the Chapter 27 tiny GPT — closing the 31-chapter loop with a measurable reward improvement on generated text.

\subsection*{Definitions}

\begin{definition}[Stochastic policy and trajectory distribution]
A parameterized stochastic policy is $\pi_\theta(a\mid s)$. A trajectory $\tau=(s_0,a_0,r_0,\ldots,s_T)$ has return $R(\tau)=\sum_{t=0}^{T-1}\gamma^t r_t$ and distribution
\[
p_\theta(\tau)=p_0(s_0)\prod_{t=0}^{T-1}\pi_\theta(a_t\mid s_t)\,P(s_{t+1}\mid s_t,a_t).
\]
The objective is $J(\theta)=\mathbb{E}_{\tau\sim p_\theta}[R(\tau)]$. The discounted state distribution $d^\pi(s)=(1-\gamma)\sum_{t\geq 0}\gamma^t \Pr_\pi[s_t=s]$ and value functions $V^\pi,Q^\pi$ are inherited from Chapter 29; the advantage is $A^\pi(s,a)=Q^\pi(s,a)-V^\pi(s)$.
\end{definition}

\begin{definition}[Score function]
The likelihood-ratio identity for the trajectory distribution is
\[
\nabla_\theta\log p_\theta(\tau)=\sum_{t=0}^{T-1}\nabla_\theta\log\pi_\theta(a_t\mid s_t),
\]
since $p_0$ and $P$ are $\theta$-independent.
\end{definition}

\begin{definition}[Generalized Advantage Estimation, Schulman 2016]
With TD residual $\delta_t=r_t+\gamma V(s_{t+1})-V(s_t)$,
\[
\hat A_t^{\mathrm{GAE}(\gamma,\lambda)}=\sum_{l=0}^{\infty}(\gamma\lambda)^l\delta_{t+l},\qquad \lambda\in[0,1].
\]
\end{definition}

\begin{definition}[TRPO and PPO surrogates]
Let $r_t(\theta)=\pi_\theta(a_t\mid s_t)/\pi_{\mathrm{old}}(a_t\mid s_t)$. TRPO maximizes $\mathbb{E}[r_t(\theta)\hat A_t]$ subject to $\mathbb{E}_s D_{\mathrm{KL}}(\pi_{\mathrm{old}}\|\pi_\theta)\leq\delta$. PPO replaces the constraint by a clipped objective:
\[
\mathcal{L}^{\mathrm{PPO}}(\theta)=\mathbb{E}\!\left[\min\!\bigl(r_t(\theta)\hat A_t,\;\mathrm{clip}(r_t(\theta),1-\varepsilon,1+\varepsilon)\hat A_t\bigr)\right].
\]
\end{definition}

\begin{definition}[GRPO, DeepSeek 2024]
For each prompt $x$ sample $G$ completions $\{y_1,\ldots,y_G\}$ from $\pi_{\mathrm{old}}(\cdot\mid x)$ with rewards $r_1,\ldots,r_G$. The group-relative advantage is
\[
\hat A_i \;=\; \frac{r_i-\mathrm{mean}(r_{1:G})}{\mathrm{std}(r_{1:G})+\epsilon},
\]
broadcast to every token of $y_i$. The PPO clipped surrogate is then applied per token, with a KL penalty $\beta\,\mathbb{E}\,D_{\mathrm{KL}}(\pi_\theta\|\pi_{\mathrm{ref}})$ to the SFT reference.
\end{definition}

\subsection*{Theorems}

\begin{theorem}[Policy gradient theorem; Sutton 1999]
\label{thm:pg}
Under conditions allowing exchange of $\nabla_\theta$ and $\int$,
\[
\nabla_\theta J(\theta)
=\mathbb{E}_{\tau\sim p_\theta}\!\left[\sum_{t=0}^{T-1}\nabla_\theta\log\pi_\theta(a_t\mid s_t)\,Q^{\pi_\theta}(s_t,a_t)\right]
=\frac{1}{1-\gamma}\,\mathbb{E}_{(s,a)\sim d^{\pi_\theta}}\!\left[\nabla_\theta\log\pi_\theta(a\mid s)\,Q^{\pi_\theta}(s,a)\right].
\]
\end{theorem}

\begin{proof}
Begin with $J(\theta)=\int p_\theta(\tau)R(\tau)\,d\tau$ and apply the log-derivative identity $\nabla p_\theta=p_\theta\nabla\log p_\theta$:
\[
\nabla_\theta J(\theta)=\int p_\theta(\tau)\,\nabla_\theta\log p_\theta(\tau)\,R(\tau)\,d\tau=\mathbb{E}[\nabla_\theta\log p_\theta(\tau)\,R(\tau)].
\]
Substitute the factorized score: $\nabla_\theta\log p_\theta(\tau)=\sum_t\nabla_\theta\log\pi_\theta(a_t\mid s_t)$ since $p_0$ and $P$ do not depend on $\theta$. Causality of rewards: $r_{t'}$ for $t'<t$ is uncorrelated with $a_t$ given history, so $\mathbb{E}[\nabla\log\pi_\theta(a_t\mid s_t)\,r_{t'}]=0$ for $t'<t$. Therefore
\[
\nabla_\theta J(\theta)=\mathbb{E}\!\left[\sum_t\nabla_\theta\log\pi_\theta(a_t\mid s_t)\sum_{t'\geq t}\gamma^{t'}r_{t'}\right]
=\mathbb{E}\!\left[\sum_t\gamma^t\nabla_\theta\log\pi_\theta(a_t\mid s_t)\,Q^{\pi_\theta}(s_t,a_t)\right],
\]
using $\mathbb{E}[\sum_{t'\geq t}\gamma^{t'-t}r_{t'}\mid s_t,a_t]=Q^\pi(s_t,a_t)$. Folding $\gamma^t$ into the discounted state distribution gives the second form.
\end{proof}

\begin{lemma}[Baselines are unbiased]
\label{lem:baseline}
For any function $b:S\to\mathbb{R}$,
\[
\mathbb{E}_{a\sim\pi_\theta(\cdot\mid s)}[\nabla_\theta\log\pi_\theta(a\mid s)\,b(s)]=0.
\]
\end{lemma}

\begin{proof}
Pull $b(s)$ outside the action expectation and use the log-trick in reverse:
\[
\sum_a\pi_\theta(a\mid s)\nabla_\theta\log\pi_\theta(a\mid s)=\sum_a\nabla_\theta\pi_\theta(a\mid s)=\nabla_\theta\sum_a\pi_\theta(a\mid s)=\nabla_\theta 1=0. \qedhere
\]
\end{proof}

Substituting $b=V^\pi$ yields the advantage form $\nabla_\theta J=\mathbb{E}[\nabla_\theta\log\pi_\theta(a\mid s)\,A^\pi(s,a)]$.

\begin{lemma}[GAE recursion and bias-variance]
$\hat A_t^{\mathrm{GAE}(\gamma,\lambda)}$ satisfies $\hat A_t=\delta_t+\gamma\lambda\,\hat A_{t+1}$. With $\lambda=0$, $\hat A_t=\delta_t$ (one-step TD; low variance, biased by bootstrap). With $\lambda=1$, telescoping gives $\hat A_t=\sum_{l\geq 0}\gamma^l r_{t+l}-V(s_t)$ (Monte Carlo minus baseline; unbiased, high variance).
\end{lemma}

\begin{proof}
Direct from $\sum_{l\geq 0}(\gamma\lambda)^l\delta_{t+l}=\delta_t+\gamma\lambda\sum_{l\geq 0}(\gamma\lambda)^l\delta_{t+1+l}$.
\end{proof}

\begin{theorem}[TRPO surrogate-improvement inequality; Schulman 2015]
For policies $\pi,\pi'$,
\[
J(\pi')\geq J(\pi)+\mathbb{E}_{(s,a)\sim d^\pi,\pi'(\cdot\mid s)}\!\left[A^\pi(s,a)\right]-\frac{2\varepsilon\gamma}{(1-\gamma)^2}\,D_{\mathrm{KL}}^{\max}(\pi,\pi'),
\]
with $\varepsilon=\max_{s,a}|A^\pi(s,a)|$ and $D_{\mathrm{KL}}^{\max}=\max_s D_{\mathrm{KL}}(\pi(\cdot\mid s)\|\pi'(\cdot\mid s))$.
\end{theorem}

\begin{proof}[Sketch]
The Kakade--Langford performance-difference lemma reads $J(\pi')-J(\pi)=\frac{1}{1-\gamma}\mathbb{E}_{(s,a)\sim d^{\pi'},\pi'}[A^\pi(s,a)]$. Replacing $d^{\pi'}$ by $d^\pi$ introduces an error bounded by the total-variation distance $\|d^{\pi'}-d^\pi\|_{TV}$. Pinsker's inequality (Chapter~11) bounds this by $\sqrt{\tfrac12 D_{\mathrm{KL}}^{\max}}$, and a discounted-chain argument supplies the $(1-\gamma)^{-2}$ factor. See Kakade--Langford (2002) and Schulman (2015) for the full proof.
\end{proof}

\begin{proposition}[PPO is a first-order trust region]
Around $\theta_{\mathrm{old}}$, $r_t(\theta)=1+\nabla_\theta\log\pi_\theta(a_t\mid s_t)\big|_{\theta_{\mathrm{old}}}^{\top}(\theta-\theta_{\mathrm{old}})+O(\|\theta-\theta_{\mathrm{old}}\|^2)$. Bounding $r_t\in[1-\varepsilon,1+\varepsilon]$ thus enforces a per-sample first-order surrogate of the KL trust region that TRPO solves with a Fisher-vector product. Empirically the clip is loose enough to be cheap and tight enough to mimic the trust region.
\end{proposition}

\begin{theorem}[GRPO group baseline is asymptotically unbiased]
Sample $y_i\sim\pi_{\mathrm{old}}(\cdot\mid x)$ i.i.d., set $\bar r=\tfrac{1}{G}\sum_j r_j$ and $\hat A_i=(r_i-\bar r)/(\hat\sigma+\epsilon)$. As $G\to\infty$, $\bar r\to\mu(x):=\mathbb{E}_{y\sim\pi_{\mathrm{old}}}[r(x,y)]$ and $\hat\sigma\to\sigma(x)$, both independent of any single $y_i$. Then
\[
\mathbb{E}_{y_i}[\hat A_i\nabla\log\pi_\theta(y_i\mid x)] = \frac{1}{\sigma(x)}\bigl(\mathbb{E}[r_i\nabla\log\pi_\theta]-\mu(x)\,\mathbb{E}[\nabla\log\pi_\theta]\bigr) = \frac{1}{\sigma(x)}\nabla J(x),
\]
since the second term vanishes by Lemma~\ref{lem:baseline}.
\end{theorem}

\begin{proof}[Remark on finite $G$]
For finite $G$, $\bar r$ depends on $y_i$ through its own term, so an $O(1/G)$ bias appears. DeepSeek (2024) shows empirically that this bias is negligible for $G\geq 4$ and that the variance reduction more than compensates. The engineering payoff is dramatic: \emph{no value network}.
\end{proof}

\begin{theorem}[DPO--RLHF equivalence; recap of Chapter 28]
The KL-regularized RLHF objective $\max_\pi \mathbb{E}_{x,y\sim\pi}[r(x,y)]-\beta D_{\mathrm{KL}}(\pi\|\pi_{\mathrm{ref}})$ has closed-form maximizer
\[
\pi^*(y\mid x)=\frac{1}{Z(x)}\pi_{\mathrm{ref}}(y\mid x)\exp\!\bigl(r(x,y)/\beta\bigr).
\]
Inverting yields $r(x,y)=\beta\log\frac{\pi^*(y\mid x)}{\pi_{\mathrm{ref}}(y\mid x)}+\beta\log Z(x)$. Substituted into the Bradley--Terry preference likelihood, $Z(x)$ cancels in the difference $r(x,y_w)-r(x,y_l)$, yielding the DPO loss directly on $\pi_\theta$ with no reward model.
\end{theorem}

\subsection*{Code sketch}

The notebook walks the algorithmic ladder. (1) REINFORCE on a 3-state MDP with softmax policy and Monte Carlo returns. (2) REINFORCE with a learned tabular state-value baseline; gradient variance drops by an order of magnitude. (3) GAE on a 5-state chain with $\lambda\in\{0,0.5,0.95,1\}$, displaying the bias-variance trade-off. (4) PPO with a 2-layer MLP policy and hand backprop on a CartPole-style env from Chapter 30. (5) GRPO warm-up on a synthetic 5-arm contextual bandit. (6) The climax: GRPO on the Chapter 27 tiny GPT, where the reward fires when the generated continuation contains a target substring; the KL to $\pi_{\mathrm{ref}}$ is monitored to confirm reward improvement without policy collapse.

\subsection*{Connection to LLMs}

Modern post-training is three coordinates on the same KL-regularized landscape:

\begin{itemize}
    \item \textbf{RLHF (PPO + RM).} SFT $\to$ reward model $\to$ PPO with KL penalty to the SFT reference. Used by InstructGPT, GPT-3.5/4, Claude, Llama-2-Chat. Pros: arbitrary, learned reward signals. Cons: four models in memory (policy, value, reward, reference); reward hacking; brittle infrastructure.
    \item \textbf{DPO.} Skip the RM; use the closed-form solution of the KL-regularized objective to write a direct preference loss on $\pi_\theta$ vs $\pi_{\mathrm{ref}}$. Used by Zephyr, T\"ulu, most open-weight chat models since late 2023. Pros: one model, one loss, no online sampling. Cons: limited to pairwise preferences; bounded by offline-dataset diversity.
    \item \textbf{GRPO.} PPO without the value network. Sample $G$ completions per prompt; baseline with the group mean; normalize by group std. Used by DeepSeek-Math and DeepSeek-R1. Pros: scales to verifiable rewards (math, code, format); no reward-model bottleneck. Cons: $G$ samples per prompt at training time.
\end{itemize}

The fellowship-level mental model: all three are surrogate optimizers of $J(\theta)-\beta D_{\mathrm{KL}}(\pi_\theta\|\pi_{\mathrm{ref}})$. RLHF estimates the gradient with PPO and a value function; DPO eliminates online sampling by inverting the closed-form solution; GRPO eliminates the value function via the group baseline. Pre-training (Chapter 27) is the floor — it sets what is sayable — and post-training reshapes only a small region around it. Pre-training is calorie intake; alignment is digestion. Every chapter in this 31-chapter chain composes to produce that final sentence.
